\section{速度と位置}\label{節:速度と位置}
今度は逆に,速度が分かっているときに位置の変化を求めてみよう.速度$v(t)=2t$で,時刻$0$から$T$までを考える.ただし$T>0$とする.速度は途中で変わるので,\,$2T\times T$では求められない.これは最初から最後まで,最後の速度で動いたと考えた計算だからである.

速度が一定なら掛け算で求められるのだから,速度があまり変わらないくらい短い時間に区切ってやればよいだろう.そこで,時間を$n$等分する.1つの区間の長さは$\frac{T}{n}$である.第$k$区間は時刻$\frac{(k-1)T}{n}$から$\frac{kT}{n}$までなので,その始めと終わりの速度はそれぞれ
\[
 \frac{2(k-1)T}{n},\qquad\frac{2kT}{n}
\]
である.速度は増加しているので,各区間を始めの速度で動き続けたと考えれば位置の変化を少なく,終わりの速度で動き続けたと考えれば多く見積もることになる.この2つの和を$L_n,U_n$と書くと,
\begin{align*}
 L_n&=\sum_{k=1}^n\frac{2(k-1)T}{n}\cdot\frac{T}{n}\\
 &=\frac{2T^2}{n^2}\cdot\frac{n(n-1)}2\\
 &=T^2\left(1-\frac1n\right),\\
 U_n&=\sum_{k=1}^n\frac{2kT}{n}\cdot\frac{T}{n}\\
 &=\frac{2T^2}{n^2}\cdot\frac{n(n+1)}2\\
 &=T^2\left(1+\frac1n\right)
\end{align*}
となる.位置の変化を$D$とすると,\,$L_n\le D\le U_n$である.\,$n$を大きくすれば両者は$T^2$に近づき,その差$U_n-L_n=\frac{2T^2}{n}$はいくらでも小さくできる.\,$D$は$n$によらない値なので,\,$D=T^2$となるはずである.最初の位置が$s(0)$なら,\,$s(T)=s(0)+T^2$となる.前節では$t^2$から速度$2t$を求めたが,逆向きの計算で再び$t^2$が現れたのだ.

この計算を図で見直してみよう.横軸を時刻,縦軸を速度とすれば,各項は長方形の「縦$\times$横」である.したがって,位置の変化を足す計算は,速度のグラフの下の面積を求める計算と同じ形になる.直線$v=2t$の下は三角形なので,その面積からも
\[
 \frac12\times T\times2T=T^2
\]
と求められる.面積の単位に対応するものは,ここでは「メートル毎秒$\times$秒」すなわちメートルである.

\begin{figure}[htbp]
\centering
\begin{tikzpicture}[x=2.7cm,y=1.25cm,>=Stealth]
 \foreach \k in {0,...,5}{
  \pgfmathsetmacro{\l}{\k/4}
  \pgfmathsetmacro{\r}{(\k+1)/4}
  \filldraw[fill=softteal,draw=tealblack!60] (\l,0) rectangle (\r,{2*\l});
 }
 \draw[->](-0.1,0)--(1.95,0)node[right]{$t$};
 \draw[->](0,-0.1)--(0,3.65)node[above]{$v$};
 \draw[very thick,tealblack](0,0)--(1.7,3.4)node[right]{$v=2t$};
 \draw[dashed](1.5,0)node[below]{$T$}--(1.5,3)--(0,3)node[left]{$2T$};
\end{tikzpicture}
\caption{各長方形は,短い時間での位置の変化を表す.}\label{図:速度の長方形}
\end{figure}

ガリレオは『新科学対話』の第3日で,静止から一様に加速する運動について,進む距離が時間の二乗に比例することを論じている.ここでの$v(t)=2t$は,係数を計算しやすく選んだ同じ種類の運動である.実際の落下をこの式で表すには,空気抵抗などを無視できる状況を考える必要がある.本書ではその物理的な妥当性を調べる代わりに,与えられた関数からどのように量を求められるかを考える.

\begin{example}{等しい時間ごとの変化}
上の運動について,時刻$k-1$から$k$までの位置の変化を求めよ.ただし$k$は正の整数とする.
\end{example}
時刻$k$までの位置の変化は$k^2$,時刻$k-1$までの変化は$(k-1)^2$なので,求める変化はその差
\[
 k^2-(k-1)^2=2k-1
\]
である.\,$k=1,2,3,\ldots$とすると$1,3,5,\ldots$となる.逆にこれらを第$n$項まで足せば,時刻$n$までの位置の変化$n^2$に戻るはずである.実際,
\begin{align*}
 \sum_{k=1}^{n}(2k-1)
 &=2\cdot\frac{n(n+1)}2-n\\
 &=n^2
\end{align*}
となる.位置の変化を引き算で求め,その変化を足して位置へ戻したのである.先ほどは時間の間隔を小さくしたが,整数の時刻だけを見れば,位置の列$0,1,4,9,\ldots$の階差が$1,3,5,\ldots$であり,階差を足すと元の列に戻るという数列の計算になっている.
